\documentclass[10pt,a4paper]{amsart}
\usepackage{amsmath,amssymb,amsthm,graphicx}
\usepackage{verbatim}
\newenvironment{reference}{\comment}{\endcomment}
%\newenvironment{reference}{}{}
\newenvironment{slogan}{\comment}{\endcomment}
\newenvironment{history}{\comment}{\endcomment}
% For commutative diagrams we use Xy-pic
\usepackage[all]{xy}

% We use 2cell for 2-commutative diagrams.
\xyoption{2cell}
% Theorem environments.
%
\theoremstyle{plain}
\newtheorem{theorem}[subsection]{Theorem}
\newtheorem{proposition}[subsection]{Proposition}
\newtheorem{lemma}[subsection]{Lemma}

\theoremstyle{definition}
\newtheorem{definition}[subsection]{Definition}
\newtheorem{example}[subsection]{Example}
\newtheorem{exercise}[subsection]{Exercise}
\newtheorem{situation}[subsection]{Situation}

\theoremstyle{remark}
\newtheorem{remark}[subsection]{Remark}
\newtheorem{remarks}[subsection]{Remarks}

\numberwithin{equation}{subsection}

% Macros
%
\def\lim{\mathop{\mathrm{lim}}\nolimits}
\def\colim{\mathop{\mathrm{colim}}\nolimits}
\def\Spec{\mathop{\mathrm{Spec}}}
\def\Hom{\mathop{\mathrm{Hom}}\nolimits}
\def\Ext{\mathop{\mathrm{Ext}}\nolimits}
\def\SheafHom{\mathop{\mathcal{H}\!\mathit{om}}\nolimits}
\def\SheafExt{\mathop{\mathcal{E}\!\mathit{xt}}\nolimits}
\def\Sch{\mathit{Sch}}
\def\Mor{\mathop{\mathrm{Mor}}\nolimits}
\def\Ob{\mathop{\mathrm{Ob}}\nolimits}
\def\Sh{\mathop{\mathit{Sh}}\nolimits}
\def\NL{\mathop{N\!L}\nolimits}
\def\CH{\mathop{\mathrm{CH}}\nolimits}
\def\proetale{{pro\text{-}\acute{e}tale}}
\def\etale{{\acute{e}tale}}
\def\QCoh{\mathit{QCoh}}
\def\Ker{\mathop{\mathrm{Ker}}}
\def\Im{\mathop{\mathrm{Im}}}
\def\Coker{\mathop{\mathrm{Coker}}}
\def\Coim{\mathop{\mathrm{Coim}}}

% Boxtimes
%
\DeclareMathSymbol{\boxtimes}{\mathbin}{AMSa}{"02}

%
% Macros for moduli stacks/spaces
%
\def\QCohstack{\mathcal{QC}\!\mathit{oh}}
\def\Cohstack{\mathcal{C}\!\mathit{oh}}
\def\Spacesstack{\mathcal{S}\!\mathit{paces}}
\def\Quotfunctor{\mathrm{Quot}}
\def\Hilbfunctor{\mathrm{Hilb}}
\def\Curvesstack{\mathcal{C}\!\mathit{urves}}
\def\Polarizedstack{\mathcal{P}\!\mathit{olarized}}
\def\Complexesstack{\mathcal{C}\!\mathit{omplexes}}
% \Pic is the operator that assigns to X its picard group, usage \Pic(X)
% \Picardstack_{X/B} denotes the Picard stack of X over B
% \Picardfunctor_{X/B} denotes the Picard functor of X over B
\def\Pic{\mathop{\mathrm{Pic}}\nolimits}
\def\Picardstack{\mathcal{P}\!\mathit{ic}}
\def\Picardfunctor{\mathrm{Pic}}
\def\Deformationcategory{\mathcal{D}\!\mathit{ef}}
\usepackage{cleveref}
\crefname{originalreference}{original reference}{original references}
\Crefname{originalreference}{Original reference}{Original references}
\makeatletter
\newlabel{algebra-lemma-universally-exact-split}{{algebra-lemma-universally-exact-split}{source}{}{sourcepointer-1}{}}
\newlabel{algebra-lemma-universally-exact-split@cref}{{[originalreference][0][]algebra-lemma-universally-exact-split}{[1][0][]source}{}{}{}}
\makeatother
\title{A map of modules whose induced map on first Tor vanishes for all test modules factors through a free module}
\author{The Stacks Project Authors}
\date{stacks-project:local-cohomology.tex:local-snapshot-20261012; GNU FDL 1.2 or later}
\begin{document}
\maketitle
\noindent\textit{Editorial note.} All delivered mathematical text is the author's own, copied from the pinned Stacks Project chapter source sources/stacks-local-cohomology/pinned-original-source.tex; no statement or proof step was rewritten, reordered or removed.\par\medskip
\noindent\textit{Editorial note.} Preamble provenance: the Stacks Project keeps one single shared preamble for every chapter; this card takes it directly from the staged sources/stacks-local-cohomology/preamble.tex as three verbatim spans (lines 8-12, 14-18 and 154-217), so no derived file was written and nothing was added or redefined. The wrapper class and the base packages are supplied by the assembler.\par\medskip
\noindent\textit{Editorial note.} Bibliography provenance: neither the counted statement nor its proof contains a citation (this lane rescans the whole span before writing the card), so the card carries no bibliography and no entry was reconstructed.\par\medskip
\noindent\textit{Editorial note.} Selection provenance: this card was opened under the file-based control contract (CONTROL.md), with the mechanical pre-card gate run on its statement and proof interval before the card was written.\par\medskip
\noindent\textit{Allowed standard prerequisites (editorial).}\begin{itemize}
\item The first Tor functor and the presentation of a finitely presented module as the cokernel of a map of finite free modules. Used at the start of the counted proof (source lines 4307-4321). Supplies the two short exact sequences that the proof compares.
\item Universal exactness of a short exact sequence and the criterion that a universally exact sequence of modules with the middle term as displayed splits. Invoked inside the counted proof at source lines 4331-4335 and attributed to a lemma of another chapter (a cross-chapter pointer). Converts the vanishing of the boundary map into a splitting, which yields the required factorisation.
\end{itemize}\medskip
\ifdefined\hypertarget\hypertarget{sourcepointer-1}{}\fi
\noindent\textit{Source reference.} Original reference algebra-lemma-universally-exact-split: sources/stacks-local-cohomology/pinned-original-source.tex, source lines 4333--4333. Navigation pointer only: the counted text refers to the Stacks item algebra-lemma-universally-exact-split, which belongs to another chapter of the project and is not part of this card. The pointer carries the target's own label as its number because this chapter's staged snapshot does not contain the other chapter's tag table; its content is not reproduced here.\par\medskip
\begin{lemma}
\label{lemma-map-tor-1-zero}
Let $R$ be a ring. Let $M \to M'$ be a map of $R$-modules with
$M$ of finite presentation
such that $\text{Tor}_1^R(M, N) \to \text{Tor}_1^R(M', N)$ is zero for all
$R$-modules $N$. Then $M \to M'$ factors through a free $R$-module.
\end{lemma}

\begin{proof}
We may choose a map of short exact sequences
$$
\xymatrix{
0 \ar[r] &
K \ar[r] \ar[d] &
R^{\oplus r} \ar[r] \ar[d] &
M \ar[r] \ar[d] &
0 \\
0 \ar[r] &
K' \ar[r] &
\bigoplus_{i \in I} R \ar[r] &
M' \ar[r] &
0
}
$$
whose right vertical arrow is the given map.
We can factor this map through the short exact sequence
\begin{equation}
\label{equation-pushout}
0 \to K' \to E \to M \to 0
\end{equation}
which is the pushout of the first short exact sequence by $K \to K'$.
By a diagram chase we see that the assumption in the lemma
implies that the boundary map $\text{Tor}_1^R(M, N) \to K' \otimes_R N$
induced by (\ref{equation-pushout}) is zero, i.e., the sequence
(\ref{equation-pushout}) is universally exact. This implies by
Algebra, Lemma \ref{algebra-lemma-universally-exact-split}
that (\ref{equation-pushout}) is split (this is where we use that
$M$ is of finite presentation). Hence the map $M \to M'$
factors through $\bigoplus_{i \in I} R$ and we win.
\end{proof}

\end{document}
